\chapter{Chờ Bạn Nhắc Mới Làm}
\chapterintrobi{Tự giác có, nhưng phải được bật công tắc}{Only act after a reminder}{Có kiểu chăm chỉ chỉ hoạt động khi nghe câu “Ê, làm bài chưa?”.}

\begin{lucbatbox}{Lục bát: Bạn nhắc mới sực nhớ}
Bài kia em vẫn để yên\
Không phải em quên, chỉ là chưa làm\
Tới khi bạn hỏi vội vàng\
À ra hôm ấy có trang phải về\
Tự nhiên ý thức tràn trề\
Nhưng sao phải đợi bạn kề bên rung
\end{lucbatbox}

\begin{tutuyetbox}{Thất ngôn tứ tuyệt: Tự giác mượn tạm}
Nếu không ai nhắc, chuyện còn xa\
Bạn vừa lên tiếng, trí em ra\
Tự giác của em không hề mất hẳn\
Chỉ là khởi động hơi qua loa
\end{tutuyetbox}

\begin{ghichu}
Chờ bạn nhắc mới làm có lợi một điểm: giúp tình bạn thêm bền. Có hại nhiều điểm hơn: làm tự chủ mãi không lớn nổi.
\end{ghichu}

\begin{englishnote}
\textbf{reminder}: lời nhắc.\par
\textbf{self-discipline}: tự giác.\par
\textbf{I forgot until my friend asked me.}: Em quên tới khi bạn em hỏi.
\end{englishnote}